% AC-G112: actual compact level and the coefficient-uniform counting interface.
% New proof. Original-K automorphy and the native counting hypotheses are retained.
\subsection{Principal compact level for the cubic tuple}
\label{sec:cubic-principal-level-counting}

The cyclotomic condition in the literal source \texttt{LevelAt} cannot
describe the ramified cubic character at 269.  The automorphic counting
space, however, is defined by a principal compact level.  We prove the
required compact-level statement directly, and identify the precise part
of the counting proof that this replaces.  This does not assert the
literal, false \texttt{LevelAt} condition at a larger integer.

Throughout this subsection let
\[
 K=\mathbb Q(\sqrt{-3}),\qquad p_0=269,\qquad
 \chi:G_K\longrightarrow\mu_3
\]
be the fixed Kummer character of \(K(\sqrt[3]{269})/K\).
The coefficient prime is denoted by \(\ell\), and \(c_\ell\) denotes
the actual cyclotomic character, of Hodge--Tate weight \(-1\).
Thus \(c_\ell(\mathrm{Frob}_v^{\rm geom})=q_v^{-1}\).
Fix \(e\geq1\), \(j\in\{1,2\}\), and \(N\geq1\).
Every rational finite character used below has conductor dividing \(N\).
In particular this applies to the original rational character \(\psi\),
and also to its prime-to-\(\ell\) Teichmuller quotient whenever that
quotient is the character actually imposed by the residual-kernel
condition.  The coefficient field may be enlarged finitely.

For a finite place \(v\) of \(K\), write
\[
 K_v(M)=\ker\bigl(\mathrm{GL}_3(\mathcal O_{K_v})
       \longrightarrow\mathrm{GL}_3(\mathcal O_{K_v}/M\mathcal O_{K_v})\bigr).
 \tag{G112.1}
\]
When \(M\) is a unit at \(v\), this means the whole maximal compact
subgroup.  The product of these groups is exactly the source
\texttt{KLev K 3 M}: its definition is the kernel over
\(\widehat{\mathcal O}_K/M\widehat{\mathcal O}_K\), mapped into the
finite adelic general linear group.

\paragraph{The actual arithmetic domain.}
We use the following established local properties of an actual
characteristic-zero cubic model point \(\rho\).  We list them to keep
the domain of the assertion independent of an unproved all-prime
extension of the model.
\begin{enumerate}
\item The representation is continuous and absolutely irreducible over
\(G_K\).  Original-\(K\) automorphy means that
\(\nu=\rho c_\ell^{-1}\) is attached, in the literal source
\texttt{AttC} convention, to a regular algebraic conjugate-self-dual
cuspidal representation \(\Pi\) of \(\mathrm{GL}_3(\mathbb A_K)\).
The exact \texttt{WitA} premise is retained when a native wrapper uses it.
Potential automorphy over a larger host alone is not this premise.
\item Away from \(\ell\), the actual residual-kernel condition gives
finite inertia, whose three characters, with multiplicities, are
\[
     \chi^j\psi_v,\quad 1,\quad\chi^j\psi_v^{-1}.
     \tag{G112.2}
\]
Here \(\psi_v\) is the restriction of the specified rational finite
character or its stated Teichmuller quotient; its conductor divides
\(N\).  These are actual inertia characters, not only their reductions.
\item At every \(v\mid\ell\), the actual representation is potentially
crystalline and has a stable full flag with graded inertia characters
\[
    \chi^j\psi_v c_\ell^e,\quad1,\quad
    \chi^j\psi_v^{-1}c_\ell^{-e}.
    \tag{G112.3}
\]
Unramified characters on the graded pieces are unrestricted.  The
parallel Hodge--Tate weights of \(\rho\) are \((-e,0,e)\).
At \(\ell=269\), these are precisely the potentially crystalline
component and full original endpoint characters in AC-G65-local/flag;
the ramified \(\chi^j\) has not been removed.  At other coefficient
primes this clause applies only on the already established local-chart
domain.  It supplies no model at the omitted primes 2 and 3.
\end{enumerate}
For the currently retained actual models, the residue characteristic is
different from 2 and 3.  The elementary assertions below also apply
whenever the three displayed characteristic-zero local properties have
been supplied by other means.

\begin{proposition}[AC-G112-local: Actual cubic principal compact level]
\label{prop:G112-local}
\textnormal{(AC-G112-local).}
In the stated arithmetic domain put \(M=269N\).  Every local
Weil--Deligne parameter of \(\nu\), after Frobenius semisimplification,
is a sum of three rank-one Weil characters and has monodromy zero.
Each corresponding smooth character of \(K_v^\times\) is trivial on
\(1+M\mathcal O_{K_v}\) when \(v\mid M\), and is unramified when
\(v\nmid M\).  Consequently
\[
                  \Pi_f^{\mathrm{KLev}(K,3,269N)}\ne0.
                  \tag{G112.4}
\]
There is no factor in this level depending on \(e\), the coefficient
field, or \(\ell\).  The exact original-\(K\) automorphy and the listed
local properties remain premises.
\end{proposition}

\begin{proof}
We first compute the finite-character conductors.  At 269 the polynomial
\(X^3-269\) is Eisenstein over the unramified quadratic extension
\(K_{269}/\mathbb Q_{269}\).  Its splitting extension over \(K_{269}\)
is cyclic, totally ramified and tame of degree three.  Therefore the
nontrivial characters \(\chi^j\) have conductor exponent one there:
they are trivial on the principal units, but not on all units.
Away from \(3\cdot269\), Kummer theory for a unit and the etaleness of
\(X^3-269\) give an unramified extension.  At 3,
\(-269=1-270\in1+27\mathbb Z_3\); the convergent expression
\[
 a=\exp\bigl(\log(-269)/3\bigr)\in1+9\mathbb Z_3
 \quad\hbox{satisfies}\quad a^3=-269.
\]
Thus \((-a)^3=269\), and \(\chi\) is actually trivial at 3.
These facts prove that the finite conductor of \(\chi\) is the unique
prime of \(K\) above 269.

Under local reciprocity, restriction of a rational character \(\psi\)
to \(K_v\) is composition with \(N_{K_v/\mathbb Q_q}\), where \(v\mid q\).
If \(q^a\mid N\) and \(x\in1+N\mathcal O_{K_v}\), then
\[
       N_{K_v/\mathbb Q_q}(x)\in1+q^a\mathbb Z_q.
       \tag{G112.5}
\]
Indeed the norm is the product of the conjugates of \(1+Ny\), each
conjugate of \(y\) being integral; subtracting one gives an element
divisible by \(N\) in \(\mathbb Z_q\).  This argument includes the
ramified place over 3.  If \(q\nmid N\), the rational character is
unramified.  Any quotient of its finite image has the same or smaller
conductor, so the Teichmuller version satisfies the same assertion.
It follows that every finite character in (G112.2)--(G112.3) is trivial
on \(1+269N\mathcal O_{K_v}\).  At 269 the exact bound is at most
\(\max(1,v_{269}(N))\); the stated \(v_{269}(269N)=v_{269}(N)+1\)
is a convenient uniform upper bound.  There is no unmentioned
coefficient-prime factor.

We next check the Weil--Deligne assertion.  Away from \(\ell\), finite
inertia gives \(N_{\rm WD}=0\).  Each inertia character in (G112.2)
is the restriction of a character of the full local Galois group:
\(\chi\) is a global character of \(G_K\), and \(\psi_v\) is rational.
Every inertia isotypic subspace is therefore Frobenius-stable, including
when characters repeat.  On each such subspace, Frobenius
semisimplification diagonalizes Frobenius.  After removing the fixed
inertia character the resulting summands are unramified rank-one
characters.  This proves the claim with multiplicities; no distinctness
of the three inertia characters has been used.

At \(v\mid\ell\), potentially crystalline representations form an exact
category with their potentially semistable period functor.  Apply that
functor over one finite extension making the representation crystalline
and use the stable flag (G112.3).  The monodromy of the whole object is
zero, and its semisimplification has the three rank-one graded objects.
The cyclotomic powers contribute only unramified Weil characters to
these rank-one objects.  A remaining graded character with trivial
inertia is unramified.  Hence its finite Weil inertia is precisely the
displayed finite character.  Twisting the whole representation by
\(c_\ell^{-1}\) likewise changes only unramified Weil factors.  This
establishes the same conductor bounds for \(\nu\).

Use full local--global compatibility.  Away from \(\ell\), this is the
usual Frobenius-semisimple compatibility for regular algebraic
conjugate-self-dual cuspidal representations.  At \(\ell\), the exact
input is Barnet-Lamb--Gee--Geraghty--Taylor, \emph{Local-global
compatibility for \(l=p\), II}, Theorem 1.1 and Corollary 3.2.
Rank three is Shin-regular by the definition in Section 1, without any
additional inequality on \(e\).  The conclusion is
\[
 \iota\,\mathrm{WD}(\nu|_{G_{K_v}})^{F\text{-ss}}
       \simeq \operatorname{rec}_{K_v}(\Pi_v\otimes|\det|^{-1}).
       \tag{G112.6}
\]
The theorem imposes no residual irreducibility, ordinary, Iwahori-fixed,
or good-level hypothesis.  Its deep proof and the underlying local
Langlands correspondence are retained imported theorems.  The
normalizing \(|\det|^{-1}\) is unramified and cannot change compact
fixed vectors or unit conductors.

Since \(\Pi\) is cuspidal on a general linear group, each \(\Pi_v\) is
generic.  The generic classification in Zelevinsky Theorem 9.7(b)
expresses it as an actually irreducible normalized induction of
unlinked essentially square-integrable segments.  Compatibility of
local Langlands with segments and parabolic induction identifies the
segment length with the monodromy block length.  Equation (G112.6)
has monodromy zero and only rank-one Weil constituents, so every
segment has length one and comes from a character of \(K_v^\times\).
Thus \(\Pi_v\) is the full irreducible principal series of three smooth
characters \(\vartheta_{v,1},\vartheta_{v,2},\vartheta_{v,3}\), whose
unit restrictions have the bounds already proved.  This does not
infer fixed vectors for an arbitrary subquotient from those of an
induced representation.

Here is the needed fixed vector.  Write \(L=K_v\), let \(B_L\) be its
upper triangular Borel, and suppose first that \(v\mid M\).  In the
normalized induced model define a section supported on \(B_LK_v(M)\)
by
\[
 f_v(bk)=\delta_{B_L}(b)^{1/2}
             \prod_{i=1}^3\vartheta_{v,i}(b_{ii})
       \quad(b\in B_L,\ k\in K_v(M)).
       \tag{G112.7}
\]
For \(b\in B_L\cap K_v(M)\), every diagonal entry belongs to
\(1+M\mathcal O_L\), so the product is one; the modular character is
one on these units as well.  Hence the formula is well-defined.
The set \(B_LK_v(M)\) is open and has compact image in \(B_L\backslash
\mathrm{GL}_3(L)\), since it is the image of the compact group
\(K_v(M)\); that image is closed in the Hausdorff quotient.  Extension
by zero is therefore a smooth induced section.  It is right
\(K_v(M)\)-invariant and \(f_v(1)=1\).
When \(v\nmid M\), all three characters are unramified, and the same
argument with \(K_v(M)=\mathrm{GL}_3(\mathcal O_L)\), using Iwasawa
decomposition, gives the usual spherical section.

The restricted tensor product of these local sections is nonzero and
fixed by the product of the groups (G112.1).  Equivalently, one can
apply the whole source \texttt{GlobFixed.globFixed} proof, which
assembles local fixed vectors and contains no \texttt{KillsMu}
condition.  This proves (G112.4).
\end{proof}

\paragraph{Exact counting inputs.}
Retain the source \texttt{UpperTheory.V.Af.Hypotheses'} and its
automorphic coefficient, cohomology, conjugation, integrality,
reciprocity and strong-multiplicity-one inputs.  We use the exact
\texttt{SigK} and \texttt{Win.EigData} definitions, and the stronger
\texttt{Af.TrData} used by the final norm theorem, with the original
arithmetic/geometric Frobenius conversion.  The separate
\texttt{Win.TrData} has a weaker size bound and is not silently
identified with \texttt{Af.TrData}.  No abstract
source hypothesis is asserted to have been proved by the local vector
construction.

\begin{proposition}[AC-G112-count: The same polynomial rank-three counting bound]
\label{prop:G112-count}
\textnormal{(AC-G112-count).}
There are constants \(C,A\), depending only on the fixed field \(K\),
rank three and the retained counting inputs, with the following
property.  Fix \(e\geq1\), \(N\geq1\), and a test prime \(q\) not
dividing \(269N\operatorname{disc}(K)\), with a place \(w\mid q\).
The actual absolutely irreducible representations in the domain of
AC-G112-local, as the coefficient prime \(\ell\ne2,3,q\) and the
coefficient field vary, admit one common source trace classifier at
\(w\), chosen before \(\ell\), with
\[
                 \#\mathcal T\leq C\{N(e+1)\}^{A}.
                 \tag{G112.8}
\]
It classifies isomorphism classes and their actual geometric-Frobenius
traces with the same integral and conjugation properties used by the
source common-integer construction.  Thus the rank-three part of that
construction survives with a fixed change of its polynomial constant.
Combining it with another trace classifier uses the unchanged source
sum/norm argument.  In particular this proposition does not silently
identify a cubic-twisted Type II endpoint with the old
\texttt{Adm2B} class.
\end{proposition}

\begin{proof}
Put \(M=269N\), \(\nu=\rho c_\ell^{-1}\), and
\[
 m=(e-1,-1,-e-1),\qquad H=e+1.
 \tag{G112.9}
\]
These are exactly the exponents in the source
\texttt{admissible\_twist\_cyclotomic} calculation.
They are strictly decreasing, satisfy \(m_i+m_{2-i}=-2\), and have
\(|m_i|\leq H\).  The actual Hodge--Tate weights of \(\nu\) are
\((1-e,1,1+e)\).  Original-\(K\) automorphy with its standard
Hodge--Tate compatibility gives the source infinitesimal character
\(\operatorname{infOf}(m)\) at every embedding.  Together with
conjugate self-duality and (G112.4), this places \(\Pi\) in the literal
set \(\mathrm{SigK}(K,m,M)\).

The complete proof of \texttt{TCardK.card\_k\_of\_printed\_results}
uses Weyl's bound for the coefficient representation, the cohomology
bound and the Borel count at \(K\mathrm{Lev}(3M)\subseteq
K\mathrm{Lev}(M)\).  Its conclusion is
\[
 \#\{\hbox{pairwise nonisomorphic }\Pi\in\mathrm{SigK}(K,m,M)\}
       \leq c\{M(H+1)\}^{A}.
       \tag{G112.10}
\]
Only the three defining properties of \texttt{SigK} enter that proof.
The exact global cyclotomic character condition \texttt{KillsMu} does
not.  Hence it applies to the space just obtained.  Moreover
\[
 c\{269N(e+2)\}^{A}
          \leq c\,538^{A}\{N(e+1)\}^{A}.
          \tag{G112.11}
\]

For clarity, we give the complete interface replacement in
\texttt{countK\_alt}.  On all finite places off \(M\), take the
algebraically normalized Satake multiset of each element of
\(\mathrm{SigK}(K,m,M)\).  The source \texttt{SatK} supplies that
multiset from the compact fixed vector.  The source \texttt{TConjK}
shows that these systems are closed under \(\operatorname{Aut}(\mathbb C)\)
at the same infinitesimal character and compact level.  It is harmless
that a complex automorphism may change an individual cubic character:
we count the entire \texttt{SigK} space, not one chosen character
idempotent.  The source \texttt{IntFinal}/\texttt{IntPHolds} supplies
the exact integrality after the source normalization; the source
\texttt{win\_algSat} supplies its size window.  Applying the same
\texttt{exists\_eigData} construction therefore gives a finite
eigenvalue data object \(D\) satisfying (G112.10), before any
coefficient prime is selected.

The only use of \texttt{LevelKNAlt} in the complete proof of
\texttt{countK\_alt} is the later local declaration \texttt{hform}:
for each Galois representation one needs a \(\Pi\in\mathrm{SigK}\),
semisimplicity, and normalized Frobenius equality off \(M\ell\).
Here original-\(K\) automorphy supplies \(\Pi\) and semisimplicity,
AC-G112-local supplies its compact level, (G112.9) supplies its
infinitesimal character, and full local--global compatibility supplies
the equality at every place off \(M\ell\).  These establish exactly
\texttt{hform}, without applying \texttt{LevelKNAlt} to an object
outside its domain.

The remainder is literal: choose this form for each representation,
map it to its system class in \(D\), and compare the systems at all
places off \(M\ell\).  Equality in \(D\) gives equality of
characteristic polynomials at a cofinite set of primes.  The source
Chebotarev/Brauer--Nesbitt theorem for continuous semisimple
representations then gives isomorphism.  At the specified test place,
\(\ell\ne q\) gives the exact Frobenius polynomial.  Thus the same
\(D\), already chosen, classifies every permitted coefficient prime.

We now construct the stronger \texttt{Af.TrData} explicitly on the
same carrier as \(D\), instead of using the weaker size field of
\texttt{Win.TrData}.  The literal \texttt{exists\_eigData} definition
uses inverse eigenvalues:
\[
  (D.\beta(\operatorname{cl}\Pi)).\operatorname{map}(b\mapsto b^{-1})
       =\operatorname{algSat}_3(\Pi_w).
       \tag{G112.12a}
\]
Set
\[
        t(s)=q_w^{-1}\sum_{b\in D.\beta(s)}b^{-1}.
        \tag{G112.12b}
\]
This is the same carrier and trace function as the source rank-three
trace construction; the strengthened norm field requires its own
proof.  Every element of this eigensystem carrier is represented by
a form in \(\mathrm{SigK}(K,m,M)\), and that form is RACSDC.  The
all-finite-place temperedness used explicitly in BLGGT I,
Proposition 1.1(2), gives three Satake roots of absolute value one
at the spherical place \(w\).  Algebraic normalization multiplies
each by \(q_w^{(3-1)/2}=q_w\).  Equations (G112.12a)--(G112.12b)
therefore express \(t(s)\) as a sum of three numbers of absolute
value one, proving \(|t(s)|\le3\).  This applies to every carrier
element, not only to the images of one coefficient-prime family.
The complete printed proof of the cited proposition was checked;
its invocation of Caraiani's temperedness theorem remains an imported
input, not a new proof of that theorem.

The \texttt{D.conj} field gives the required conjugation permutation
for \(t\), since automorphisms fix \(q_w\) and commute with inversion.
For integrality use the literal \texttt{D.intInv} field at \(a=e-1\):
every \(m_i\le e-1\), so each \(q_w^{e-1}b^{-1}\) is integral.
Their sum is \(q_w^e t(s)\).  Together with the finite carrier,
these are exactly all fields of \texttt{Af.TrData} in
\texttt{Classifies.lean}.  No field is borrowed from the weaker
\texttt{Win.TrData} without this verification.

If
\(\sigma\) is the arithmetic Frobenius used by source \texttt{IsFrob},
then \(c_\ell(\sigma)=q_w\).  Consequently
\[
 \operatorname{tr}\nu(\sigma^{-1})
      =q_w\operatorname{tr}\rho(\sigma^{-1}).
      \tag{G112.12}
\]
The division by \(q_w\) is exactly the one in (G112.12b);
there is no new Tate or inverse-Frobenius convention.  Its carrier has
the same cardinality as \(D\).  Equations (G112.10)--(G112.12) give
(G112.8) with \(C=c538^A\), enlarging constants if the retained
source count already used a harmless factor.

The final algebra in \texttt{KeepNumU} forms a sum of supplied
trace classifiers.  Here we use the literal \texttt{Af.TrData.sum}
and apply the \texttt{Af.TrData} version of
\texttt{frobenius\_trace\_norm\_nonzero\_log\_bound} directly.  This latter
argument knows only the trace-data axioms, their cardinalities, the
test-prime size conditions and the chosen finite Galois-stable root
set.  Our rank-three data satisfy precisely those axioms.  Hence its
integer is again selected before \(\ell\), and its logarithmic bound
has the same polynomial form after (G112.11).  If other classes are
included, their actual trace-data hypotheses must be independently
supplied; no old rank-two admissibility statement is altered here.
\end{proof}

\paragraph{Source and application boundary.}
The source \texttt{LevelAt} requires finite characters factoring through
the action on rational roots of unity.  At 269 there is an inertia
element fixing all roots of unity and the rational finite characters
but having nontrivial cubic value.  Thus increasing its positive
integer \(M\) never repairs that definition.  AC-G112 replaces its
specific use by (G112.4); it does not contradict the obstruction.
The source \texttt{Hypotheses'}, original-\(K\) automorphy and its
normalization, and the actual local-chart domain remain explicit.
No all-family or all-prime \(M4\), reducedness, integral Selmer-length
bound, omitted-prime estimate, or final counting budget follows from
this subsection.  The fixed integer 269 changes the counting constant,
not the polynomial exponent; the test primes must also avoid 269.

\paragraph{Geometric dictionary.}
The finite inertia characters describe the ramification of the
associated local systems.  Principal level records a finite
congruence cover on which an automorphic vector is fixed; it is not
the assertion that this ramification comes from a cyclotomic cover.
Only the already audited finite-cover category is being used in this
description.  No new theorem of Grothendieck is attributed here.
